\section{Homology localizers and truncated homotopy types}

\begin{claim}{C.1}{Known}{The families and their model structures}
For a $\mathcal U$-small spectrum $E$, we use the unreduced convention
\[
 \W_E=\{u:A\to B\mid E\wedge\Sigma^\infty(BA)_+
       \longrightarrow E\wedge\Sigma^\infty(BB)_+\text{ is an equivalence}\}.
\]
For a unital ring $R$, not assumed commutative and including the zero ring, $\W_R$ means
$\W_{HR}$. For an integer $n\geq0$, $\W_n$ consists of the functors inducing
a bijection on $\pi_0$ and isomorphisms on all based $\pi_i$, $1\leq i\leq n$.
These are accessible fundamental localizers. Their associated simplicial
localizations have monomorphisms as cofibrations. On $\Cat$ they admit the
left proper Thomason model structure, with the usual Thomason cofibrations.
For $\W_n$ this structure is also right proper.
\end{claim}
The existence assertions are \citet[Theorem 9.4.2, Corollary 9.4.3,
Theorem 5.2.15]{RootCisinski2006}. For the truncations, the defining sphere
generator and properness are stated in \citet[9.2.1, Corollary 9.2.15,
Proposition 9.2.17]{RootCisinski2006}. In particular the sphere generator
is $\Delta/\partial\Delta^{n+2}$, of homotopy type $S^{n+1}$; the shift
in its index is essential. We use these existence theorems with their
proofs in the cited source, rather than re-proving the small-object
argument. The remaining axiom assertions below have explicit proofs.

\begin{claim}{C.2}{Proved}{Sums, connectedness, filtered colimits, and accessibility}
Each $\W_n$, $n\geq0$, and each $\W_E$ satisfies Loc~(4), (8), and (9).
Each $\W_n$ satisfies Loc~(5). For $\W_E$, Loc~(5) holds exactly when
$E\not\simeq0$. Thus it holds for $\W_R$ exactly when $R\neq0$.
\end{claim}
\begin{proof}
The nerve of a disjoint union is the disjoint union of the nerves.
Generalized homology of a disjoint union is the direct sum of the
unreduced homology groups. A block-diagonal map of direct sums is an
isomorphism if and only if each block is. This proves both implications
of Loc~(4), including arbitrary small index sets and the empty family.
The same assertion for $\W_n$ follows by taking components and based
homotopy groups component by component.

Suppose $E\neq0$ and choose $k\in\mathbb Z$ with $\pi_kE\neq0$.
If a map of spaces $f:X\to Y$ omits a component $Y_j$, the summand
$E_k(Y_j)$ contains a copy of $\pi_kE$, split by any point and the
augmentation, and cannot be hit by $E_k(f)$. If distinct components
$X_i,X_{i'}$ map into the same component of $Y$, choose one point in
each. Their images are joined by a path. For nonzero $v\in\pi_kE$,
the two point inclusions define the nonzero element $(v,-v)$ in the
corresponding direct-sum summands of $E_k(X)$, and its image is zero.
Thus an $E$-equivalence is bijective on components. This also proves
that $\varnothing\to X$ is an $E$-equivalence only when $X=\varnothing$.
For $E=0$ every functor belongs to $\W_E$, and the fold $e\amalg e\to e$
violates Loc~(5).

The $q$-simplices of a nerve are finite strings of composable arrows.
Consequently $N$ commutes with filtered colimits. Filtered colimits of
simplicial sets model their filtered homotopy colimits: every finite
simplicial sphere, disk, and homotopy factors through one stage, and
equalities between finitely many simplices hold at a later stage.
This proves commutation of components and based homotopy groups with
the filtered colimit and establishes Loc~(8) for $\W_n$. Suspension
spectra and smashing with $E$ preserve homotopy colimits, and stable
homotopy groups commute with filtered homotopy colimits. This proves
Loc~(8) for $\W_E$. Loc~(9) is the accessibility assertion of
\claimref{C.1}; here accessibility is for the full subcategory of
$\Cat^{[1]}$, as in XVI, manuscript pp.~74--78.
\end{proof}

\begin{claim}{C.3}{Proved}{The cofibration columns}
Every $\W_E$, including $E=0$, satisfies Loc~(6 a), (6 b), and (6 c).
Every $\W_n$, $n\geq0$, satisfies Loc~(6 a) and fails both Loc~(6 b)
and Loc~(6 c).
\end{claim}
\begin{proof}
For a covering by two sieves $X_1,X_2\subset X_3$, with intersection
$X_0$, every simplex of $NX_3$ lies in one of $NX_1,NX_2$. Indeed the
last object of its chain lies in one sieve, and closure under sources
then places every earlier object in that sieve. The nerve square is
therefore a pushout along simplicial monomorphisms, hence a homotopy
pushout. After applying $E\wedge\Sigma^\infty(-)_+$ it yields the
cofiber sequence of spectra
\[
 E\wedge\Sigma^\infty(BX_0)_+
 \longrightarrow
 (E\wedge\Sigma^\infty(BX_1)_+)\oplus
 (E\wedge\Sigma^\infty(BX_2)_+)
 \longrightarrow E\wedge\Sigma^\infty(BX_3)_+.
\]
A map of these sequences satisfies two out of three for equivalences.
Equivalence of the middle direct-sum map is equivalent to equivalence
of its two diagonal blocks. These statements prove precisely the three
implications (6 a), (6 b), and (6 c). No connectivity or nonemptiness
assumption was made.

For $\W_n$, homotopy pushouts preserve $n$-equivalences because
$\tau_{\leq n}:\cS\to\cS_{\leq n}$ is a left adjoint. This proves
(6 a) on the same nerve squares. To disprove the other two clauses,
let $P_n$ be the poset of nonempty proper subsets of
$\{0,\ldots,n+1\}$, ordered by inclusion. Its objects, arrows, and
composition are thus completely specified: there is one arrow for
each inclusion and composites are inclusions. Its nerve is the
barycentric subdivision of $\partial\Delta^{n+1}$ and realizes to
$S^n$. Adjoin two distinct maximal objects $+$ and $-$, each above
every object of $P_n$, with no arrow between them. Denote the resulting
poset by $S(P_n)$ and its two cones by $C_+(P_n),C_-(P_n)$.
These cones are sieves, cover $S(P_n)$, have intersection $P_n$,
and have terminal objects. Their union realizes to the unreduced
suspension $S^{n+1}$.

Map this cover to the cover $e=e\cup e$. The maps from both cones and
from $S(P_n)$ to $e$ are $n$-equivalences, whereas $P_n\to e$ is not.
This is the full cube counterexample to (6 b). For (6 c), map the
cover $C_-(P_n)=P_n\cup C_-(P_n)$ into
$S(P_n)=C_+(P_n)\cup C_-(P_n)$. The intersection map is $\id_{P_n}$,
the union map $C_-(P_n)\hookrightarrow S(P_n)$ is an $n$-equivalence,
and the first piece map $P_n\hookrightarrow C_+(P_n)$ is not.
The construction includes $n=0$, where $P_0$ has two objects and
$S(P_0)$ realizes to a circle. The program
\texttt{code/ce\_sphere\_cover.py} checks the finite category laws,
sieve conditions, and nerves of the first four instances.
\end{proof}

\begin{claim}{C.4}{Proved}{A nullity obstruction valid for every nonzero spectrum}
Let $E\neq0$, and let $A$ be an $E$-aspherical space. Then $A$ is
nonempty and connected, and
$\Hom(\pi_1(A,a),\mathbb Z)=0$ for every $a\in A$. For a free group
$F$, the constant-map inclusion
$BF\to\Map(A,BF)$ is an equivalence. Thus $BF$ is null with respect
to every $E$-aspherical space.
\end{claim}
\begin{proof}
Connectedness and nonemptiness follow from \claimref{C.2}. If a
homomorphism $\pi_1(A,a)\to\mathbb Z$ is nonzero, its image is
$m\mathbb Z$ for some $m\geq1$. Dividing by $m$ gives a surjection
$\chi:\pi_1(A,a)\twoheadrightarrow\mathbb Z$. Choose a loop mapping
to $1$. The corresponding pointed maps
$s:S^1\to A$ and $h:A\to K(\mathbb Z,1)=S^1$ satisfy
$hs\simeq\id_{S^1}$, since pointed maps to $K(\mathbb Z,1)$ are
classified by their homomorphisms on fundamental groups. Therefore
$\Sigma E$ is a retract of $E\wedge\Sigma^\infty A$, where the latter
uses the chosen basepoint. This last spectrum is zero by
$E$-asphericity, a contradiction.

Any nontrivial subgroup $H$ of a free group is free: it is the
fundamental group of the connected covering graph associated with
$H$. A nontrivial free group admits a surjection to $\mathbb Z$.
Consequently the image of any homomorphism $\pi_1(A,a)\to F$ is
trivial. Mapping to the $1$-type $BF$ depends only on the fundamental
groupoid of $A$. Since $A$ is connected, this groupoid is
pointedly equivalent to the one-object groupoid
$B\pi_1(A,a)$. The pointed mapping groupoid then has the
homomorphisms $\pi_1(A,a)\to F$ as objects and conjugacies
with conjugating element forced to be $1$ as morphisms.
Thus $\Map_*(A,BF)$ is the discrete space of these
homomorphisms, and is therefore contractible.
Evaluation at $a$ has this mapping space as
fiber, so the constant-map section is an equivalence.
\end{proof}

\begin{claim}{C.5}{Proved}{No nonzero homology localizer has enough intrinsic fibrations}
For every spectrum $E\neq0$, $\W_E$ fails Loc~(7 F). In particular
this holds for $\W_R$ for every nonzero unital ring $R$,
with no restriction on its characteristic or torsion.
\end{claim}
\begin{proof}
Let $F=\langle x,y\rangle$ and let
\[
 \alpha:F\longrightarrow F,\qquad
 \alpha(x)=x,\quad \alpha(y)=yxyx^{-1}y^{-1}.
\]
The completely specified group categories, the proof that
$B\alpha$ is an integral homology equivalence, and the finite
$S_3$ witness proving that $\alpha$ is not surjective are given in
\claimref{F.8}. A graph model $BF\simeq S^1\vee S^1$ shows directly
that its reduced suspension-spectrum map is the identity matrix:
the exponent-sum matrix of $\alpha$ is $I_2$. Hence $B\alpha$ is an
$E$-equivalence for every spectrum, including nonconnective spectra.
By \claimref{C.4}, both copies of $BF$ are local for nullification
at all $E$-aspherical objects. This nullification exists and is
accessible by \citet[Proposition 6.1.8]{RootCisinski2006}, since
$\W_E$ is accessible. A localization is fully faithful on
its local objects, so it cannot invert this non-equivalence between
them. Thus $\W_E$ is not generated by its aspherical categories.
The equivalence of aspherical generation and Loc~(7 F), proved in
\claimref{E.2}, completes the argument.
\end{proof}

\begin{claim}{C.6}{Proved}{Failure of the stronger fibration axioms}
For $E\neq0$, and for every finite $n\geq0$, both Loc~(7) and the
weakest Loc~(7 $D_{\Winfty}$) fail for $\W_E$ and $\W_n$.
For $E=0$ all the fibration clauses hold.
\end{claim}
\begin{proof}
Here is a direct counterexample to the weakest clause for $\W_E$.
With $F$ and $\alpha$ as in \claimref{C.5}, define a groupoid
$\mathcal P$ with object set $F$, and with
$\Hom_{\mathcal P}(g,g')=F$. Write $h:g\to g'$ for the arrow labelled
by $h$. Its composite with $h':g'\to g''$ is labelled $h'h$.
Define
\[
 p:\mathcal P\longrightarrow BF,\qquad
 p(h:g\to g')=g'\alpha(h)g^{-1}.
\]
This is a functor because the formula respects multiplication.
The functor $j:BF\to\mathcal P$ at object $1$, sending $h$ to its
label, is an equivalence of categories, and $pj=B\alpha$.
Thus $p\in\W_E$. For an arrow $k$ in the base, the label-$1$
arrow $g\to kg$ lifts it. These lifts and their inverses make $p$
both a Grothendieck fibration and opfibration, with isomorphic
transport functors; in particular $p\in\Parf_{\Winfty}$.
Its strict fiber over the unique object has components
$F/\alpha(F)$, since its arrows satisfy
$g'\alpha(h)g^{-1}=1$. There are at least two components, because
$\alpha$ is not surjective. This fiber is not $E$-aspherical by
\claimref{C.2}. Hence $p$ is in $\Parf_{\Winfty}\cap\W_E$ and is
not universally in $\W_E$, exactly the failure of
Loc~(7 $D_{\Winfty}$).

For $\W_n$, the map $*\to K(\mathbb Z,n+1)$ is an $n$-equivalence,
but its homotopy fiber $K(\mathbb Z,n)$ is not $n$-aspherical
(for $n=0$ this fiber is the discrete set $\mathbb Z$).
Represent the map by a functor between simplex categories and
apply the $\Winfty$-equivalence--$\Parf_{\Winfty}$ factorization
described in \claimref{E.4}. The resulting perfect functor remains
in $\W_n$, its strict fibers model those homotopy fibers, and so it
violates the same clause. Loc~(7) implies this clause by Part A.
When $E=0$, all functors belong to $\W_E$, every functor is in
$\Fib_{\W_E}$, and the identity factorization proves (7 F).
\end{proof}

\begin{claim}{C.7}{Proved}{The complete axiom table for the named families}
With the conventions of \claimref{C.1}, the following table is exact.
In its $7D$ column we use the weakest version
$7D_{\Winfty}$; failure there also means failure of all stronger
$7D$ versions.
\[
\begin{array}{c|cccccccccc}
 &4&5&6a&6b&6c&7&7D&7F&8&9\\\hline
\W_R,\ R\ne0&+&+&+&+&+&-&-&-&+&+\\
\W_E,\ E\ne0&+&+&+&+&+&-&-&-&+&+\\
\W_n,\ 0\leq n<\infty&+&+&+&-&-&-&-&+&+&+\\
\W_0&+&+&+&-&-&-&-&+&+&+\\
\W_E,\ E=0&+&-&+&+&+&+&+&+&+&+
\end{array}
\]
\end{claim}
\begin{proof}
All entries except $7F$ for $\W_n$ are proved in
\claimref{C.1}--\claimref{C.6}. The sphere generator of
\claimref{C.1} makes $\W_n$ aspherically generated, so
\claimref{E.2} proves that remaining entry. Repeating the $\W_0$
row records explicitly the case singled out in XVI, manuscript
pp.~68--70 and 91. In particular the answer to the manuscript's
question about $\W_n$ and (7 F) is affirmative, while the answer
for every nonzero homology theory is negative.
\end{proof}

\begin{claim}{C.8}{Proved}{An intrinsic description of $\Fib_{\W_0}$}
Let $p:T\to B$ be a functor. For a component $B_j\subset B$, put
$T_j=T\times_BB_j$, and for $b\in B_j$ let $T_b$ be the strict fiber.
Then $p\in\Fib_{\W_0}$ if and only if both of the following conditions
hold. First, $\pi_0(T_b)\to\pi_0(T_j)$ is a bijection for every $b$.
Second, for every arrow $\beta:b\to b'$ and every
$c\in\pi_0(T_j)$, there exists an arrow of $T$ over $\beta$ whose
source, and hence target, lies in $c$.
Equivalently, for every functor $X\to B$ the canonical map is a bijection
\[
 \pi_0(T\times_BX)\longrightarrow
 \pi_0(T)\times_{\pi_0(B)}\pi_0(X).
\]
\end{claim}
\begin{proof}
Suppose $p\in\Fib_{\W_0}$. Pulling back the $\W_0$-map
$e_b\to B_j$ proves the first condition. Pulling back either
endpoint inclusion $e\to[1]$, where $[1]\to B$ represents $\beta$,
shows that both fibers induce bijections on components into
$T\times_B[1]$. Its only arrows between its two layers are the
arrows of $T$ over $\beta$. The first condition labels the fiber
components by $c\in\pi_0(T_j)$; therefore each such label must have
an arrow between the layers. This proves the second condition.

Conversely, fix a component $X_a$ mapping into $B_j$. For each
$x\in X_a$, the objects over $x$ labelled by $c$ form a nonempty
connected category. An arrow $x\to x'$ in $X_a$ has a lift joining
these two categories by the second condition. A finite zigzag in
$X_a$ therefore joins all objects with a given label $c$ in
$T\times_BX_a$. Distinct labels cannot be joined, since their
images in $T_j$ have distinct components. This proves the displayed
bijection, naturally in $X\to B$. A $\W_0$-equivalence over $B$
induces a bijection on its right-hand side, and hence on its
left-hand side. Empty fibers and empty components obey the same
argument: if $T_j$ is empty, all the relevant fibers are empty and
both conditions are vacuous there.
\end{proof}

\begin{claim}{C.9}{Proved}{The $n$-equivalence fibration class}
Let $n\geq0$ and $p:T\to B$ be a functor. Then
$p\in\Fib_{\W_n}$ if and only if, for every $X\to B$, the canonical
comparison is an equivalence in $\cS_{\leq n}$:
\[
 \tau_{\leq n}B(T\times_BX)\longrightarrow
 \tau_{\leq n}BT
 \mathbin{\times_{\tau_{\leq n}BB}}
 \tau_{\leq n}BX.                                      \tag{C.9.1}
\]
The right-hand pullback is a homotopy pullback of $n$-types.
In particular every $\W_n$-Thomason fibration belongs to
$\Fib_{\W_n}$. This supplies a functorial factorization for (7 F)
at every finite $n$.
\end{claim}
\begin{proof}
Right properness in \claimref{C.1} implies that a strict pullback
along a model fibration preserves all weak equivalences. It also
implies that such strict pullbacks represent derived pullbacks,
by taking a fibrant replacement of the base and comparing the two
pullback diagrams by right properness. Thus Thomason fibrations
have the asserted property.

For $p\in\Fib_{\W_n}$, factor $p$ over $B$ as a $\W_n$-equivalence
$T\to\widehat T$ followed by a Thomason fibration
$\widehat T\to B$. Both $T$ and $\widehat T$ are objects of the
Brown fibration category over $B$ with intrinsic fibrations
$\Fib_{\W_n}$. The dual Brown lemma of Part A says that a weak
equivalence between two such fibrations is preserved by every
strict base change. Consequently
$T\times_BX\to\widehat T\times_BX$ is a $\W_n$-equivalence.
The preceding model-category calculation for $\widehat T$ proves
(C.9.1), including its canonical map. Conversely, if (C.9.1) is
an equivalence for every $X$, apply it to the source and target of
any $\W_n$-map over $B$. Their right-hand sides are equivalent,
so the strict pullback map is a $\W_n$-map. This is precisely the
definition of $\Fib_{\W_n}$.
This is the specialization to $\W_n$ of the intrinsic
weak-fibration criterion in
\citet[Proposition 6.2.3]{RootCisinski2006}; its properness
hypothesis is supplied by \claimref{C.1}.
\end{proof}

\begin{claim}{C.10}{Proved}{Bounds for the homological fibration classes}
For any spectrum $E$, $\Fib_{\W_E}$ contains all product projections
$B\times K\to B$, all inclusions of unions of components, and the
closure of these maps under composition, strict base change, and
retracts. More generally it contains every functor $p$ for which
$Np$ universally preserves $E$-equivalences of simplicial sets
under strict pullback. An upper bound is the class $Q_E$ of
functors $p:T\to B$ such that, for every $m\geq0$, every
$\sigma:[m]\to B$, and every $i\in[m]$, the strict fiber inclusion
\[
 T_{\sigma(i)}\longrightarrow T\times_B[m]
\]
belongs to $\W_E$: one has $\Fib_{\W_E}\subseteq Q_E$.
For every $E\neq0$ this upper inclusion is strict.
Indeed $\Fib_{\W_E}$ omits the explicit
three-sheeted action-groupoid covering described below, although
that covering lies in $\Parf_{\Winfty}\cap Q_E$.
For a fixed $p$, membership in $\Fib_{\W_E}$ can be tested on a set
of bounded-cardinality $\W_E$-maps over its base, rather than on a
proper class.
\end{claim}
\begin{proof}
For a product projection the pullback map is $u\times\id_K$.
The suspension spectrum of its nerve, after adjoining a basepoint,
is $\Sigma^\infty(Bu)_+\wedge\Sigma^\infty(BK)_+$. Smashing an
$E$-equivalence with the second factor preserves equivalence.
Component inclusions follow from Loc~(4), and the three closure
properties follow directly from the definition of $\Fib$ and
closure of $\W_E$ under retracts. The nerve preserves strict
pullbacks, proving the simplicial sufficient condition. It is
often called a sharpness condition; we do not identify it with
the unrelated assertion that a map is merely a fibration in a
left Bousfield model structure that is not right proper.
For the upper bound, the functor $[0]\xrightarrow{i}[m]$
over $B$ is in $\Winfty\subseteq\W_E$, so its pullback along
any $p\in\Fib_{\W_E}$ has to be in $\W_E$. This is exactly
the defining condition of $Q_E$, with all its quantifiers.

Let $F=\langle x,y\rangle$, $S=\{1,2,3\}$, and let
$\rho:F\to S_3$ send $x$ to $(12)$ and $y$ to $(13)$.
Define the action groupoid $S/\!/F$ to have objects $S$ and
arrows labelled $g\in F$ from $s$ to $\rho(g)s$, with composition
given by multiplication in $F$. Its projection
$p:S/\!/F\to BF$ is a discrete fibration and opfibration,
with bijective transports. The action is transitive, so its
total groupoid is connected. Pulling back along $B\alpha$ from
\claimref{C.5} replaces the action by $\rho\alpha$, which sends
both generators to $(12)$. Its orbits are $\{1,2\}$ and $\{3\}$.
Thus the pullback map is not a $\W_0$-map, hence not a $\W_E$-map
for $E\neq0$, although $B\alpha\in\W_E$. The group laws and
these orbit calculations are checked in
\texttt{code/f\_free\_group\_witness\_output.json}.
To check that this same covering lies in $Q_E$, pull it back
along $\sigma:[m]\to BF$. Transport from vertex $0$
identifies the resulting category with the disjoint union of
three copies of $[m]$. Every strict fiber inclusion is
therefore a weak homotopy equivalence. Thus the upper bound
is strictly larger than the intrinsic fibration class.

Finally fix $B$. Accessibility of $\W_E\subset\Cat^{[1]}$
gives a regular cardinal $\kappa$ and a set $\mathcal T$ of
$\W_E$-maps whose $\kappa$-filtered colimits exhaust $\W_E$;
choose $\kappa$ so that the inclusion preserves these colimits.
The cardinalities of the objects and arrows in all sources and
targets of members of $\mathcal T$ have a common bound, since
$\mathcal T$ is a set. Equip each member of $\mathcal T$ with
every possible functor from its target to $B$; these choices
again form a set. For a given $u$ over $B$, any chosen
presentation by $\mathcal T$ becomes a presentation over $B$
by composing the stage-to-target maps with the specified map
to $B$. Strict pullback along $p$ commutes with these filtered
colimits, because objects and arrows are computed by finite
limits of sets. Loc~(8) therefore reduces the test to that set
of bounded stages over $B$. This is an exact set-sized test,
not a claimed finite decision procedure.
\end{proof}

\begin{claim}{C.11}{Proved}{The cofibration threshold at $\W_1$}
Every functor is a $\W_0$-cofibration. If a fundamental localizer
$\W$ contains $\W_1$, every functor injective on objects is a
$\W$-cofibration. Hence Loc~(6 A) holds in this range. These facts
do not license replacing the products with endpoint inclusions
in XII, manuscript pp.~228--234, by pushout-products.
\end{claim}
\begin{proof}
The functor $\pi_0:\Cat\to\mathbf{Set}$ is left adjoint to the
discrete-category inclusion and preserves all pushouts. Pushout
in sets preserves an isomorphism along every map. This proves
the first assertion.

For the second, groupoid completion $\Pi_1:\Cat\to\mathbf{Gpd}$
preserves pushouts and does not change objects. The folk model
structure on groupoids has equivalences as weak equivalences,
object injections as cofibrations, and all objects cofibrant and
fibrant. Its pushout along an object injection computes the
homotopy pushout. One can see the relevant gluing assertion
directly by representing morphisms in the groupoid pushout as
alternating words and cancelling adjacent arrows in each
factor. Consequently, for a span in $\Cat$ with one leg
injective on objects, the canonical comparison from its
Grothendieck homotopy colimit to its ordinary pushout is a
$\W_1$-equivalence. The comparison is therefore in $\W$.
Homotopy colimits preserve objectwise $\W$-equivalences, and
two out of three now shows that the given object injection
preserves every $\W$-map under cobase change. A categorical
mapping cylinder provides a factorization by an object
injection followed by a categorical homotopy equivalence,
proving (6 A). The last sentence is a statement about the
specific construction: none of these assertions states that
the pushout-product of two designated maps is a cofibration
for a localizer below $\W_1$.
\end{proof}
